% Realization Calculus: mathematical monograph visual system.
% This file styles the presentation. The mathematical source files remain separate.
\definecolor{RCPlum}{HTML}{3B283E}
\definecolor{RCMulberry}{HTML}{76526B}
\definecolor{RCCopper}{HTML}{AB8564}
\definecolor{RCInk}{HTML}{2F2B31}
\definecolor{RCMuted}{HTML}{69646B}
\definecolor{RCCream}{HTML}{F8F6F2}
\definecolor{RCLine}{HTML}{DFD9D4}
\definecolor{RCPaper}{HTML}{FFFFFF}
\colorlet{ink}{RCInk}
\colorlet{navy}{RCPlum}
\colorlet{accent}{RCMulberry}
\colorlet{muted}{RCMuted}
\colorlet{KomorebiInk}{RCInk}
\colorlet{KomorebiPine}{RCPlum}
\colorlet{KomorebiMuted}{RCMuted}
\geometry{letterpaper,left=.84in,right=.84in,top=.89in,bottom=.82in,headheight=17pt,headsep=17pt,footskip=24pt}
\setlength{\parindent}{0pt}
\setlength{\parskip}{6.1pt plus 1pt minus 1pt}
\linespread{1.08}
\setlength{\emergencystretch}{3em}
\setlength{\headheight}{17pt}
\pagestyle{fancy}
\fancyhf{}
\fancyhead[L]{\sffamily\fontsize{8}{10}\selectfont\color{RCMuted}\MakeUppercase{Realization Calculus}}
\fancyhead[R]{\sffamily\fontsize{8}{10}\selectfont\color{RCMuted}Special Functions}
\fancyfoot[L]{\sffamily\fontsize{8}{10}\selectfont\color{RCMuted}Function Families and Their Structure}
\fancyfoot[R]{\sffamily\fontsize{8}{10}\selectfont\color{RCPlum}\thepage}
\renewcommand{\headrulewidth}{.24pt}
\renewcommand{\footrulewidth}{0pt}
\fancypagestyle{plain}{%
 \fancyhf{}\fancyfoot[R]{\sffamily\small\color{RCPlum}\thepage}%
 \renewcommand{\headrulewidth}{0pt}%
}
\titleformat{\section}{\rmfamily\Large\color{RCPlum}}{\thesection}{.7em}{}
\titleformat{\subsection}{\rmfamily\large\bfseries\color{RCInk}}{\thesubsection}{.7em}{}
\titleformat{\subsubsection}{\rmfamily\normalsize\bfseries\color{RCInk}}{\thesubsubsection}{.7em}{}
\titlespacing*{\section}{0pt}{2.1ex plus .8ex minus .2ex}{.9ex}
\titlespacing*{\subsection}{0pt}{1.8ex plus .6ex minus .2ex}{.6ex}
\hypersetup{colorlinks=true,linkcolor=RCPlum,citecolor=RCPlum,urlcolor=RCMulberry,
  pdftitle={Function Families and Their Structure},
  pdfauthor={Komorebi P.},
  pdfsubject={Realization Calculus: Special Functions}}

\renewcommand{\KomorebiCover}{%
 \clearpage\thispagestyle{empty}\begingroup
 \begin{tikzpicture}[remember picture,overlay]
 \fill[RCPlum] (current page.south west) rectangle (current page.north east);
 \draw[RCCopper,line width=1.1pt] ([xshift=.82in,yshift=-.87in]current page.north west) -- ([xshift=2.03in,yshift=-.87in]current page.north west);
 \draw[RCCopper!50,line width=.45pt] ([xshift=.82in,yshift=.72in]current page.south west) -- ([xshift=-.82in,yshift=.72in]current page.south east);
 \foreach \k in {0,...,4}{
   \draw[RCCopper!23,line width=.55pt] ([xshift=-.75in,yshift=-1.20in-\k*.24in]current page.north east)
      .. controls ++(-.6in,-.1in) and ++(.3in,.4in) .. ([xshift=-2.4in,yshift=-1.80in-\k*.16in]current page.north east);
 }
 \end{tikzpicture}
 \vspace*{.32in}
 {\sffamily\fontsize{8.5}{13}\selectfont\color{RCCopper}\bfseries REALIZATION CALCULUS\par}
 \vspace*{1.58in}
 {\color{white}\fontsize{35}{41}\selectfont\bfseries Function Families\\[.11em] and Their Structure\par}
 \vspace{19pt}
 {\color{white!78!RCCream}\fontsize{14}{19}\selectfont Coupled growth, arithmetic depth,\\ finite resolution\par}
 \vspace{29pt}{\color{RCCopper}\rule{1.18in}{.75pt}}\par
 \vspace{13pt}{\color{white}\sffamily\fontsize{9}{12}\selectfont MATHEMATICAL MONOGRAPH / SPECIAL FUNCTIONS\par}
 \vfill
 {\color{white}\fontsize{12}{16}\selectfont Komorebi P.\par}
 \vspace{.47in}
 \endgroup\clearpage
}

\newcommand{\MonographTitleLeaf}{%
 \thispagestyle{empty}\begingroup
 \vspace*{1.03in}
 {\sffamily\fontsize{8}{12}\selectfont\color{RCMulberry}\bfseries SPECIAL FUNCTIONS\par}
 \vspace{.50in}
 {\color{RCPlum}\fontsize{31}{36}\selectfont Function Families\\ and Their Structure\par}
 \vspace{.21in}
 {\color{RCMuted}\fontsize{14}{19}\selectfont Coupled growth, arithmetic depth,\\ and finite resolution\par}
 \vspace{.37in}{\color{RCCopper}\rule{.91in}{.7pt}}\par
 \vfill
 {\fontsize{12}{16}\selectfont Komorebi P.\par}
 \vspace{8pt}{\color{RCMuted}\sffamily\fontsize{9}{13}\selectfont A research monograph in five parts\par}
 \vspace{.84in}
 \endgroup\clearpage
}

\newcommand{\MonographPreface}{%
 \phantomsection\addcontentsline{toc}{section}{Preface}
 {\sffamily\fontsize{8}{11}\selectfont\color{RCMulberry}\bfseries PREFACE}\par\vspace{6pt}
 {\fontsize{25}{30}\selectfont\color{RCPlum}Five different families, one problem of structure\par}
 \vspace{14pt}
 The theory of special functions is built from examples whose definitions carry more information than their initial values suggest. A contraction hidden in a differential equation can determine the growth of an entire function. A prime omitted from a sum can alter its parameter signature. An apparently positive finite moment calculation can conceal a negative mass elsewhere in the underlying arithmetic measure.

 This volume studies these phenomena through explicit families. Some results are extensions of familiar analytic constructions. Others are concrete applications of classical machinery in new coordinates. The proofs keep those distinctions visible. A useful representation does not by itself establish historical novelty, and a numerical illustration does not substitute for a universal argument.

 The exposition begins with scalar nonlinear feedback, passes through complex arithmetic depth and signed prime measures, and reaches spectral zero orbits. It then develops coupled functions and precise limits on finite observation. A final part collects the longer proofs, reproducible calculations, and connections to existing literature.

 \vspace{10pt}
 {\color{RCPlum}\bfseries A note on evidence}\par
 Analytic identities and their domains are stated explicitly. Finite computations are reported with their actual ranges, and unresolved regularity, identifiability, and priority questions remain open. These qualifications are part of the mathematics rather than temporary editorial notes.

 \vspace{10pt}
 {\color{RCPlum}\bfseries Reading this book}\par
 The main chapters are designed to be read consecutively. The technical appendices can be consulted at the point where a proof needs more detail. References and verification records are collected at the back so that the development itself reads as a continuous monograph.
 \clearpage
}

\renewcommand{\pageheading}[2]{%
 \clearpage\phantomsection\addcontentsline{toc}{section}{#2}\markright{#2}
 {\sffamily\fontsize{8}{11}\selectfont\color{RCMulberry}\bfseries\MakeUppercase{#1}\par}
 \vspace{6pt}
 {\rmfamily\fontsize{20.5}{26}\selectfont\color{RCPlum}#2\par}
 \vspace{7pt}{\color{RCCopper}\rule{.87in}{.65pt}}\par
 \vspace{9pt}
}
\renewcommand{\claim}[1]{\par\vspace{2.5pt}\noindent{\color{RCMulberry}\bfseries #1}\par\vspace{1pt}}
\newcommand{\MonographPart}[3]{%
 \clearpage\phantomsection\addcontentsline{toc}{part}{Part #1: #2}\thispagestyle{empty}
 \vspace*{1.45in}
 {\sffamily\fontsize{9}{12}\selectfont\color{RCMulberry}\bfseries PART #1\par}
 \vspace{16pt}
 {\rmfamily\fontsize{30}{35}\selectfont\color{RCPlum}#2\par}
 \vspace{18pt}{\color{RCCopper}\rule{1.05in}{.8pt}}\par
 \vspace{17pt}
 {\fontsize{12}{18}\selectfont\color{RCMuted}#3\par}
 \vfill
 {\sffamily\fontsize{8}{12}\selectfont\color{RCMuted}REALIZATION CALCULUS\enspace /\enspace SPECIAL FUNCTIONS\par}
 \clearpage
}
\newcommand{\MonographSubpart}[1]{%
 \clearpage\phantomsection\addcontentsline{toc}{section}{#1}
 {\sffamily\fontsize{9}{12}\selectfont\color{RCMulberry}\bfseries REFERENCE SECTION\par}
 \vspace{8pt}{\fontsize{24}{28}\selectfont\color{RCPlum}#1\par}
 \vspace{12pt}{\color{RCCopper}\rule{1in}{.7pt}}\par\vspace{14pt}
}
